\chapter*{Abstract}

In today's competitive world, stress significantly impacts human health and can lead to conditions such as diabetes, depression, anxiety, and cardiovascular diseases. Timely stress recognition is thus vital for effective management. While wearable devices are commonly used to assess stress via physiological signals, they are often costly and require physical contact, causing discomfort over time. This work introduces a tabular deep learning architecture that detects stress by analyzing physiological features extracted from remote photoplethysmography (rPPG) signals derived from facial videos. The method is evaluated on the UBFC-Phys dataset for two experiments: (i) stress task classification and (ii) multi-level stress classification. In both cases, it outperforms current state-of-the-art approaches.

Additionally, rPPG has emerged as a promising contactless biometric modality by enabling cardiac signal extraction from facial videos. Although effective for physiological monitoring, its biometric application in real-world conditions is still underexplored. We propose a CNN-based embedding model trained with a triplet loss function to learn discriminative feature vectors from heartbeat-level rPPG segments. The model is evaluated on PURE, UBFC-Phys, and MSPM datasets, encompassing motion, lighting variation, and stress. It achieves near-perfect results on PURE and strong cross-session performance. However, stress-induced and unconstrained scenarios show performance degradation due to signal variability. These results demonstrate rPPG’s biometric potential and highlight challenges related to emotional and physiological states.

\vspace{1em}
\noindent\textit{\textbf{Keywords:} Stress detection, Remote photoplethysmography (rPPG), Physiological features, Biometric Identification, Signal quantization, Triplet loss.}







% In today's competitive world, stress is a major factor negatively impacting human health. Over time, it can lead to serious health issues such as diabetes, depression, anxiety, and cardiovascular diseases. Hence, timely stress recognition is crucial for effective stress management. Currently, various wearable devices are employed for stress assessment by capturing physiological signals. Although accurate, these devices are often expensive and require direct physical contact, which may lead to discomfort over extended periods. In this work, we propose a tabular-based deep learning architecture for detecting stress by analyzing physiological features. These features are extracted from remote photoplethysmography (rPPG) signals computed from facial videos. The architecture is validated on the publicly available UBFC-Phys dataset through two experiments: (i) stress task classification and (ii) multi-level stress classification. In both experiments, our proposed methodology outperforms current state-of-the-art methods.

% Remote photoplethysmography (rPPG) has also emerged as a promising contactless biometric modality by enabling the extraction of cardiac signals from facial videos. While rPPG has shown effectiveness in physiological monitoring, its application in biometric identification under real-world conditions remains underexplored. In this work, we propose a CNN-based embedding model trained with a triplet loss function to extract discriminative feature vectors from heartbeat-level rPPG segments. Our approach is evaluated on three publicly available datasets—PURE, UBFC-Phys, and MSPM—covering diverse conditions including motion, lighting variations, and psychological stress. The model achieves near-perfect results on the PURE dataset, demonstrating strong intra- and cross-session identification performance, even under quantization. However, performance degrades in stress-inducing and unconstrained scenarios due to increased signal variability. These findings confirm the potential of rPPG signals for biometric identification in controlled settings and highlight the challenges posed by emotional and physiological variability, encouraging future research into hybrid and stress-resilient biometric systems.
% \vspace{1em}
% \noindent\textit{\textbf{Keywords:} Stress detection, Remote photoplethysmography (rPPG), Physiological features, Biometric Identification, Signal quantization, Triplet loss.}

% Keywords:Remote photoplethysmography (rPPG) \sep Biometric identification \sep Signal quantization \sep Triplet loss.
% \vspace*{0.180cm}
% \noindent
% {\large{\textbf{\emph{Keywords:}}}}~  CNN, rPPG.
 








